% =============================================================================
% IJCB 2026 MBAge Challenge — 1st Place Methodology Submission
% Karthik Sivarama Krishnan & Koushik Sivarama Krishnan
% One-page methodology document. Compile on Overleaf with pdflatex.
% =============================================================================
\documentclass[9pt,letterpaper]{extarticle}

% --- Layout ---------------------------------------------------------------
\usepackage[letterpaper,top=0.4in,bottom=0.4in,left=0.55in,right=0.55in]{geometry}
\usepackage[parfill]{parskip}
\setlength{\parindent}{0pt}
\setlength{\parskip}{2pt plus 1pt minus 0.5pt}

% --- Typography -----------------------------------------------------------
\usepackage[T1]{fontenc}
\usepackage[utf8]{inputenc}
\usepackage{helvet}
\renewcommand{\familydefault}{\sfdefault}
\usepackage{microtype}

% --- Math & symbols -------------------------------------------------------
\usepackage{amsmath}
\usepackage{amssymb}
\usepackage{bm}

% --- Tables, color, links -------------------------------------------------
\usepackage{booktabs}
\usepackage{array}
\usepackage{xcolor}
\usepackage[colorlinks=true,urlcolor=blue!60!black,linkcolor=blue!60!black]{hyperref}
\usepackage{colortbl}

% --- Section headings -----------------------------------------------------
\usepackage[compact]{titlesec}
\titleformat{\section}{\small\bfseries\color{navyblue}}{}{0pt}{}
\titlespacing*{\section}{0pt}{2pt plus 0.5pt}{0pt}

% --- Color palette --------------------------------------------------------
\definecolor{navyblue}{HTML}{1A3461}
\definecolor{tablehead}{HTML}{1A3461}
\definecolor{tablerow}{HTML}{F4F6F9}
\definecolor{tablehighlight}{HTML}{FFF3CD}
\definecolor{codebg}{HTML}{F3F4F7}
\definecolor{captiongray}{HTML}{3A4452}

% --- Custom command for code-style boxed math ----------------------------
\newcommand{\codebox}[1]{%
  \begin{center}
  \colorbox{codebg}{\parbox{0.92\linewidth}{\centering\small #1}}
  \end{center}}

% --- Suppress page numbers (single-page document) ------------------------
\pagestyle{empty}

% =============================================================================
\begin{document}

% --- Title block ----------------------------------------------------------
\begin{center}
{\normalsize\bfseries\color{navyblue}
A Hierarchical Multimodal Ensemble with FiLM-Conditioned Fusion\\
for Tongue and Face Age Estimation}\\[2pt]
{\footnotesize\color{captiongray}
1st Place Solution \,\textperiodcentered\, IJCB 2026 MBAge Challenge
(CodaBench Competition \#15116)}\\[4pt]
{\small\bfseries
Karthik Sivarama Krishnan$^{1}$ \,\textperiodcentered\, Koushik Sivarama Krishnan$^{2}$}\\[1pt]
{\footnotesize\itshape
$^{1,2}$Independent Researcher \,|\, contact via competition email}
\end{center}

% --- Abstract -------------------------------------------------------------
\section*{Abstract}
We present the winning solution to the IJCB 2026 MBAge challenge, in which we
achieved first place on the final combined leaderboard ($S_{\mathrm{final}}=0.4129$)
with the top Phase~2 score of $S_{2}=0.4396$. Our approach treats age estimation
as a structured prediction problem and combines three complementary deep models
trained on pre-extracted tongue (four CNN backbones) and face (Gabor magnitude
banks) features: a hierarchical decade-then-fine Deep Label Distribution Learning
(DLDL) model, a feature-engineered dual-backbone model incorporating fourteen
statistical and attention-based descriptors, and a Feature-wise Linear Modulation
(FiLM) ensemble that median-aggregates six seed-varied face-conditioned
predictors. A weighted ensemble of the three components, with weights refined
by a 114-candidate grid search constrained to out-of-fold non-regression, yielded
our final submission. Code, predictions, and reproducibility artifacts are
publicly released.

% --- Method ---------------------------------------------------------------
\section*{Method}
\textbf{Hierarchical decade-routed DLDL (40\% weight).}
The first component decomposes regression into a five-way decade classification
(20s--60s) followed by ten-bin within-decade Deep Label Distribution
Learning~[1]. A shared three-layer MLP trunk (256 hidden units,
GELU, LayerNorm, dropout 0.3) processes global-average-pooled features from four
CNN backbones (ResNet18, MobileNetV2, EfficientNetB0, DenseNet121) concatenated
with Gabor face features (cheek and chin banks B/C/D, each 40-D). Predictions
follow soft hierarchical routing,
\[
\hat{y} \;=\; \sum_{d=1}^{5} P(d)\!\left(\mathrm{start}_d + \sum_{k=0}^{9} P(k\mid d)\cdot k\right).
\]
Training combines decade cross-entropy, within-decade KL divergence against
Gaussian soft labels ($\sigma$ annealed $2.0\!\rightarrow\!0.5$), $L_1$
regression, and Wing loss~[2], optimized with AdamW under
ten-fold stratified cross-validation (200 epochs, patience 30).

\textbf{Feature-engineered dual-backbone model (42.5\% weight).}
The second component is a DenseNet121 + ResNet18 dual-backbone hierarchical
model (hidden dim 384, dropout 0.35) augmented with all fourteen of the
following: pre-pool LayerNorm; standard-deviation and $2\!\times\!2$ grid
pooling; discrete-cosine and Haar wavelet projections; higher-order and spatial
moments; Gabor-bank marginals and cross-bank asymmetry; iterative-square-root
covariance pooling~[3] (reduce-dim 32); Compact Bilinear
Pooling~[4] (dim 1024); mutual-information channel selection
(top-64); Efficient Channel Attention~[5]; CKA-based backbone
de-duplication; learned-query attention pooling; confidence-weighted late fusion;
and the Multi-Modal Transfer Module (MMTM)~[6] for cross-backbone
fusion.

\textbf{FiLM-conditioned median ensemble (17.5\% weight).}
The third component is a seven-model median over six Feature-wise Linear
Modulation~[7] transformers
($d_{\text{model}}\!=\!128$, four heads, two layers, triple-head decoder
combining DLDL, regression, and ordinal predictions) trained with seeds
$\{0,1,2,3,4,42\}$, plus the hierarchical model. FiLM projects face features
into per-channel affine parameters $\gamma,\beta$ that modulate the tongue
representation. Median aggregation is robust to single-seed outliers and reduces
prediction variance relative to arithmetic averaging.

\textbf{Final blend.}
Component predictions were combined as a weighted average and clipped to the
valid age range:
\codebox{$\displaystyle \hat{y}\;=\;\mathrm{clip}\bigl(0.400\,\hat{y}_{\text{hier}}
+ 0.425\,\hat{y}_{\text{s12}} + 0.175\,\hat{y}_{\text{film6}},\;\,20,\;\,69\bigr).$}
Weights were selected by exhaustive grid search at 0.025 resolution over the
ternary simplex (114 candidate combinations), retaining the smallest deviation
from the reference (40/40/20) weighting that yielded a positive out-of-fold
jittered $S$ delta with no negative per-metric component. The refinement
contributed $+0.00022$ to the final leaderboard score over the equal-weight
reference, validating the proven blend structure as a tight local optimum.

% --- Results --------------------------------------------------------------
\section*{Results}
\begin{table}[h!]
\vspace{-4pt}
\centering
\footnotesize
\renewcommand{\arraystretch}{1.05}
\begin{tabular}{@{}lcccccc@{}}
\toprule
\rowcolor{tablehead}
\textcolor{white}{\textbf{Configuration}} &
\textcolor{white}{\textbf{Rank}} &
\textcolor{white}{\textbf{S}} &
\textcolor{white}{\textbf{MSE}} &
\textcolor{white}{\textbf{1Y-ACC}} &
\textcolor{white}{\textbf{5Y-ACC}} &
\textcolor{white}{\textbf{10Y-ACC}} \\
\midrule
\rowcolor{tablerow}
Phase 1 (tongue only)        & 2nd & 0.3727 & 97.17 & 0.16 & 0.53 & 0.74 \\
Phase 2 (tongue + face)      & \textbf{1st} & \textbf{0.4396} & \textbf{46.92} & \textbf{0.176} & \textbf{0.648} & \textbf{0.864} \\
\rowcolor{tablehighlight}
\textbf{Combined} ($0.4\,S_1 + 0.6\,S_2$) & \textbf{1st} & \textbf{0.4129} & --- & --- & --- & --- \\
\bottomrule
\end{tabular}\\[3pt]
{\footnotesize\color{captiongray}\textbf{Table~1.} Final CodaBench leaderboard
metrics. The competition score is computed as
$S = 0.1/(1+\mathrm{MSE}) + 0.4\,\mathrm{ACC}_{1\mathrm{Y}}
+ 0.3\,\mathrm{ACC}_{5\mathrm{Y}} + 0.2\,\mathrm{ACC}_{10\mathrm{Y}}$.
Our team (ks7585) won Phase 2 outright and the combined ranking by a margin
of $+0.0027$ over the second-place team.}
\end{table}

% --- Reproducibility ------------------------------------------------------
\section*{Reproducibility and Code Availability}
All eight component models were trained on a single NVIDIA RTX 5070 (12\,GB)
using PyTorch 2.11.0+cu128 and Python 3.12, in a total wall-clock time of
approximately 3.5\,hours. Reproducibility was verified by retraining the
hierarchical model from scratch on different hardware and observing
byte-identical out-of-fold predictions to the original training run
(MD5 \texttt{9ee818c7a9557ea0852ff9634bd62227}), confirming full determinism
of the recipe under fixed random seeds. The complete code base, including all
training scripts, out-of-fold predictions, the final submission file, and
step-by-step reproduction instructions, is released at
\href{https://github.com/ks7585/MBAge-IJCB2026}{\nolinkurl{github.com/ks7585/MBAge-IJCB2026}}
(MIT License).

% --- References -----------------------------------------------------------
\section*{Key References}
\begin{footnotesize}
\setlength{\parskip}{0pt}
\noindent
[1]~B.-B.~Gao \emph{et al.}, ``Deep Label Distribution Learning with Label
Ambiguity,'' \emph{IEEE Trans. Image Processing}, 26(6):2825--2838, 2017.\;
[2]~Z.-H.~Feng \emph{et al.}, ``Wing Loss for Robust Facial Landmark
Localisation,'' \emph{CVPR}, 2018.\;
[3]~P.~Li \emph{et al.}, ``Towards Faster Training of Global Covariance Pooling
Networks by Iterative Matrix Square Root Normalization,'' \emph{CVPR}, 2018.\;
[4]~Y.~Gao \emph{et al.}, ``Compact Bilinear Pooling,'' \emph{CVPR}, 2016.\;
[5]~Q.~Wang \emph{et al.}, ``ECA-Net: Efficient Channel Attention for Deep
Convolutional Neural Networks,'' \emph{CVPR}, 2020.\;
[6]~H.~R.~V.~Joze \emph{et al.}, ``MMTM: Multimodal Transfer Module for CNN
Fusion,'' \emph{CVPR}, 2020.\;
[7]~E.~Perez \emph{et al.}, ``FiLM: Visual Reasoning with a General Conditioning
Layer,'' \emph{AAAI}, 2018.
\end{footnotesize}

\end{document}
